\subsection{DUAL OF A MODULE OBTAINED BY EXTENSION OF SCALARS}

Let A, B be two rings, \(\rho: \mathbf{A} \to \mathbf{B}\) a ring homomorphism, E a left A-module and \(\mathbf{E}^*\) its dual. We shall define a canonical B-linear mapping

\[
\nu _ {E}: (E ^ {*}) _ {(B)} \rightarrow (E _ {(B)}) ^ {*}.\tag{19}
\]

The left hand side of (19) may be written as \(\mathrm{Hom}_{\mathbf{A}}(\mathbf{E},\mathbf{A})\otimes_{\mathbf{A}}\rho_{*}(\mathbf{B}_{s})\) , where, in \(\mathrm{Hom}_{\mathbf{A}}(\mathbf{E},\mathbf{A})\) , \(\mathbf{A}\) is considered as an \((\mathbf{A},\mathbf{A})\) -bimodule. Then there is a canonical \(\mathbf{Z}\) -homomorphism ( \(\S 4\) , no. 2, formula (7))

\[
\nu: \operatorname{Hom} _ {A} (E, A) \otimes_ {A} \rho_ {*} (B _ {s}) \rightarrow \operatorname{Hom} _ {A} (E, A \otimes_ {A} \rho_ {*} (B _ {s})) = \operatorname{Hom} _ {A} (E, \rho_ {*} (B _ {s}))
\]

with the identification given by the canonical isomorphism of § 3, no. 4, Proposition 4. On the other hand, the right hand side of (19) may be written as \(\mathrm{Hom}_{\mathbf{B}}(\rho_{*}(\mathbf{B}_{d}) \otimes_{\mathbf{A}} \mathbf{E}, \mathbf{B}_{s})\) ; as B is a (B, A)-bimodule, there is a canonical Z-isomorphism (§ 4, no. 1, Proposition 1)

\[
\beta : \operatorname{Hom} _ {\mathbf {B}} (\rho_ {*} (\mathbf {B} _ {d}) \otimes_ {\mathbf {A}} \mathbf {E}, \mathbf {B} _ {s}) \rightarrow \operatorname{Hom} _ {\mathbf {A}} (\mathbf {E}, \operatorname{Hom} _ {\mathbf {B}} (\mathbf {B} _ {s}, \mathbf {B} _ {s}))
\]

and \(\mathrm{Hom}_{\mathbb{B}}(\mathbf{B}_s,\mathbf{B}_s)\) is canonically identified, as an A-module, with \(\rho_{*}(\mathbf{B}_{s})\) (see the proof of no. 1, Proposition 1). Taking account of these identifications, we obtain the homomorphism \(\nu_{\mathbb{E}}\) ; it is easily verified that this homomorphism is characterized by the equation

\[
\langle \xi \otimes x, \nu _ {E} (x ^ {*} \otimes \eta) \rangle = \xi_ {\rho} (\langle x, x ^ {*} \rangle) \eta ,\tag{20}
\]

for \(x \in \mathbf{E}\) , \(x^{*} \in \mathbf{E}^{*}\) , \(\xi\) , \(\eta\) in \(\mathbf{B}\) , which shows immediately that \(\nu_{\mathbf{E}}\) is B-linear. Moreover, for every A-linear mapping \(u: \mathbf{E} \to \mathbf{F}\) the diagram

\[
\begin{array}{c c c} (\mathrm{F} ^ {*}) _ {(\mathrm{B})} & \xrightarrow {\nu_ {\mathrm{F}}} & (\mathrm{F} _ {(\mathrm{B})}) ^ {*} \\ \Big \downarrow^ {(t _ {u}) _ {(\mathrm{B})}} & & \Big \downarrow^ {t (u _ {(\mathrm{B})})} \\ (\mathrm{E} ^ {*}) _ {(\mathrm{B})} & \xrightarrow {\nu_ {\mathrm{E}}} & (\mathrm{E} _ {(\mathrm{B})}) ^ {*} \end{array}
\]

is commutative.

PROPOSITION 8. If one of the A-modules E, \(\rho_{*}(B_{s})\) is projective and finitely generated, the homomorphism \(\nu_{\mathbb{E}}\) is bijective.

This follows from the above and § 4, no. 2, Proposition 2.

Suppose in particular that E is a finitely generated free A-module and let \((e_{i})_{1\leqslant i\leqslant n}\) be a basis of E and \((e_{i}^{*})\) the dual basis; then the canonical isomorphism (19) maps the basis \((e_{i}^{*}\otimes1)\) of \((\mathrm{E}^{*})_{(\mathrm{B})}\) to the dual basis of the basis \((1\otimes e_{i})\) of \(\mathrm{E}_{(\mathrm{B})}\) .

